Fix a block and condition on its realized history.  This proof does not invoke \(\texttt{def:model-class}\) or \(\texttt{thm:weighted-gaussian-transport}\).  Write the two coordinates specified by \(\texttt{def:adaptive-witness}\) as \(a_{gb}=(a_{Dg},a_{Sg})^\top\), and write \(Y_{gb}(d,q):=Y_{gbi}(d,q)\), which is independent of \(i\) in this construction.  Direct substitution into the four contrast coefficients gives
\[
 m_{gbq_H}-m_{gbq_L}=(a_{Dg},a_{Sg})^\top.                 \tag{22}
\]
Indeed, the first coordinate is
\[
 \frac12\{Y_{gb}(1,q_H)-Y_{gb}(0,q_H)
             -Y_{gb}(1,q_L)+Y_{gb}(0,q_L)\}
 =\frac{\theta_g(q_H-q_L)}2=a_{Dg},
\]
and the second is
\[
 \frac12\sum_{q\in\{q_L,q_H\}}\{Y_{gb}(0,q)+Y_{gb}(1,q)\}
 =2h_g+\delta_g+\frac{\theta_g}{2}=a_{Sg}.
\]
The outcomes are constant over units in a group.  Since the second stage fixes exactly \(np_{1q}=nq\) treated and \(np_{0q}=n(1-q)\) untreated units, for each \(d\),
\[
 \frac1n\sum_i\frac{\mathbf1\{A_{gbi}=d\}Y_{gbi}(d,q)}{p_{dq}}
 =Y_{gbi}(d,q).
\]
Thus \(\widehat m_{gbq}=m_{gbq}\), \(e_{gbq}=0\), \(U_b=0\), and \(R_b=A_b\).

At either saturation, \(q=1/2+x_{gb}/4\) and \(q^2-q=-3/16\).  Substitution in the group mean gives
\[
 \bar Y^{\mathrm{obs}}_{gb}=\frac{a_{Sg}}2+\frac{a_{Dg}}4+\frac L4x_{gb}.
\]
The four type means and every balanced-sign mean are zero.  Since A and B blocks alternate and \(\lambda=1\), the predictable score is therefore
\[
 V_{gb}^{A}=\frac32H_g+\frac14J_g+x_{g,b-1},
 \qquad V_{gb}^{B}=\frac54H_g+\frac12J_g,
\]
apart from the initial score \(V_{g1}=H_g\).

We now prove the required uniform transport.  Put
\[
 y_g=\left(a_{Dg},a_{Sg},\frac{V_{gb}}{\sigma_b}\right)^\top,
 \qquad S_G=G^{-1/2}\sum_gx_{gb}y_g=(A_b^\top,W_b)^\top.
\]
For both block types, \(\sum_gy_g=0\).  The first two coordinates have absolute value at most one, and the construction gives \(|V_{gb}|\leq11/4\).  For an A block, put \(d_g=3H_g/2+J_g/4\).  Equal representation of the four types gives \(G^{-1}\sum_gd_g^2=37/16\), while the preceding assignment has Euclidean norm \(\sqrt G\).  Hence the reverse triangle inequality gives
\[
 \sigma_b=\frac{\|d+x_{b-1}\|_2}{\sqrt{G-1}}
 \geq\sqrt{\frac G{G-1}}\left(\frac{\sqrt{37}}4-1\right)>0.
\]
For a B block, direct calculation gives \(\sigma_b^2=29G/\{16(G-1)\}\), and the initial block has \(\sigma_1^2=G/(G-1)\).  These bounds are uniform over every possible history.  Consequently there is a constant \(K<\infty\), independent of \(G\), the block, and the history, such that \(\max_g\|y_g\|\leq K\).

Represent the uniform balanced-sign law as follows.  Draw a uniform perfect matching \(\mathfrak m\) of \(\{1,\ldots,G\}\), label its pairs \(\{r_j,s_j\}\), and independently draw Rademacher orientations \(\epsilon_j\), setting \(x_{r_j}=\epsilon_j\) and \(x_{s_j}=-\epsilon_j\).  Every balanced sign vector is produced by exactly \((G/2)!\) oppositely signed matchings and one orientation for each such matching.  Since there are \((G-1)!!\,2^{G/2}\) matching-orientation pairs, every member of \(\mathcal X_G\) has probability
\[
 \frac{(G/2)!}{(G-1)!!\,2^{G/2}}
 =\binom{G}{G/2}^{-1}.
\]
Conditional on \(\mathfrak m\), therefore,
\[
 S_G=\sum_{j=1}^{G/2}\epsilon_jd_j,
 \qquad d_j=G^{-1/2}(y_{r_j}-y_{s_j}),
 \qquad C_{\mathfrak m}=\sum_jd_jd_j^\top.                  \tag{23}
\]
Let \(C_G=\operatorname{Cov}_0(S_G)\).  A fixed unordered edge belongs to a uniform perfect matching with probability \(1/(G-1)\); consequently the identity
\[
 \sum_{g<h}(y_g-y_h)(y_g-y_h)^\top
 =G\sum_g(y_g-\bar y)(y_g-\bar y)^\top
\]
gives
\[
 \mathbb E_{\mathfrak m}C_{\mathfrak m}
 =\frac1{G-1}\sum_g(y_g-\bar y)(y_g-\bar y)^\top=C_G.      \tag{24}
\]
The matching covariance also concentrates uniformly.  To see this entry by entry, write it as \(G^{-1}\sum_{g<h}I_{gh}c_{gh}\), where \(|c_{gh}|\leq4K^2\).  For distinct edges, the joint matching probabilities are zero when they share a vertex and \(1/\{(G-1)(G-3)\}\) when they are disjoint.  Thus the diagonal edge terms number \(O(G^2)\) and have variance \(O(G^{-1})\), the overlapping ordered pairs number \(O(G^3)\) and have covariance \(O(G^{-2})\), and the disjoint ordered pairs number \(O(G^4)\) and have covariance
\[
 \frac1{(G-1)(G-3)}-\frac1{(G-1)^2}=O(G^{-3}).
\]
After the factor \(G^{-2}\), all three contributions are \(O(G^{-1})\).  Summing the nine entry bounds yields
\[
 \mathbb E_{\mathfrak m}\|C_{\mathfrak m}-C_G\|_F^2
 \leq \frac{C K^4}{G}.                                    \tag{25}
\]

Let \(Z_{\mathfrak m}=\sum_jg_jd_j\), where the \(g_j\)'s are independent standard normals.  For any \(h:\mathbb R^3\to\mathbb R\) with bounded derivatives through order three, replace the summands in (23) one at a time.  The constant, linear, and quadratic Taylor terms agree because \(\mathbb E\epsilon_j=\mathbb Eg_j=0\) and \(\mathbb E\epsilon_j^2=\mathbb Eg_j^2=1\).  Taylor's integral remainder and \(\|d_j\|\leq2K/\sqrt G\) give
\[
 \left|\mathbb E\{h(S_G)-h(Z_{\mathfrak m})\mid\mathfrak m\}\right|
 \leq C\|D^3h\|_\infty\sum_j\|d_j\|^3
 \leq \frac{C K^3\|D^3h\|_\infty}{\sqrt G}.              \tag{26}
\]
Let \(Z_G\sim N_3(0,C_G)\).  For two positive-semidefinite covariance matrices \(C_0,C_1\), differentiation along \(C_t=(1-t)C_0+tC_1\), followed by one Gaussian integration by parts, gives
\[
 \frac{d}{dt}\mathbb Eh\{N(0,C_t)\}
 =\frac12\operatorname{tr}\left[(C_1-C_0)
                  \mathbb E D^2h\{N(0,C_t)\}\right].       \tag{27}
\]
For singular matrices, apply the same density calculation to \(C_t+\delta I_3\) and let \(\delta\downarrow0\); boundedness of \(D^2h\) justifies dominated convergence.  Equations (25)--(27) and Cauchy--Schwarz therefore imply
\[
 \left|\mathbb E_0h(S_G)-\mathbb Eh(Z_G)\right|
 \leq\frac{C_K}{\sqrt G}
       \{\|D^2h\|_\infty+\|D^3h\|_\infty\}.              \tag{28}
\]
This comparison is on the covariance range and nowhere divides by a score variance.

Taking \(h(s)=f_{\eta,\kappa}(s_3)\) in (28) proves the scalar weighted comparison because this fixed smooth function has bounded derivatives of every order.  For the quadratic comparison, choose a smooth cutoff \(\chi_L\) equal to one on \(\|s_{1:2}\|\leq L\) and zero on \(\|s_{1:2}\|\geq2L\), and apply (28) to
\[
 h_{rs,L}(s)=s_rs_s f_{\eta,\kappa}(s_3)\chi_L(s_{1:2}),
 \qquad r,s\in\{1,2\}.
\]
The derivatives are bounded for each fixed \(L\).  Conditional on a matching, the independent-sign fourth-moment identity gives
\[
 \mathbb E\left(\sum_j\epsilon_jd_{jr}\right)^4
 =3\left(\sum_jd_{jr}^2\right)^2-2\sum_jd_{jr}^4\leq12K^4.
\]
The covariance bound following from (23) gives the same uniform fourth-moment bound for \(Z_G\).  Hence, for both laws,
\[
 \mathbb E\{\|S_{G,1:2}\|^2\mathbf1(\|S_{G,1:2}\|>L)\}
 +\mathbb E\{\|Z_{G,1:2}\|^2\mathbf1(\|Z_{G,1:2}\|>L)\}
 \leq C_KL^{-2}.                                            \tag{29}
\]
First let \(G\to\infty\) in (28) with \(L\) fixed and then let \(L\to\infty\) in (29).  Equivalently, choosing a sufficiently slowly increasing deterministic \(L=L_G\) produces one deterministic vanishing modulus.  This proves the two asserted weighted comparisons uniformly over blocks and histories, including singular score covariance.

It remains to turn the comparison into the block covariance expansion.  Complementation \(x\mapsto-x\) preserves \(P_0\), leaves the even weight \(f_{\eta,\kappa}(W_b)\) unchanged, and changes \(A_b\) to \(-A_b\).  Since \(R_b=A_b\), the soft-law score is exactly centered and
\[
 \Sigma_{b,\mathrm{RR}}
 =\frac{\mathbb E_0\{A_bA_b^\top f_{\eta,\kappa}(W_b)\}}
        {\mathbb E_0f_{\eta,\kappa}(W_b)}.                 \tag{30}
\]
The Gaussian pair has \(W_b^\star\sim N(0,1)\) and admits, on its covariance range,
\[
 A_b^\star=\beta_bW_b^\star+E_b^\star,
 \qquad E_b^\star\perp W_b^\star,
 \qquad\operatorname{Var}(E_b^\star)=\Sigma_{b,0}-\beta_b\beta_b^\top.
\]
For \(Z\sim N(0,1)\), differentiation of \(\mathbb Ee^{-\kappa Z^2}=(1+2\kappa)^{-1/2}\) gives
\[
 \mathbb E\{Z^2e^{-\kappa Z^2}\}=(1+2\kappa)^{-3/2}.
\]
Therefore the Gaussian version of (30) equals
\[
 \Sigma_{b,0}-(1-v_{\eta,\kappa})\beta_b\beta_b^\top.       \tag{31}
\]
Both the exact and Gaussian denominators are at least \(\eta\).  The just-proved numerator and denominator comparisons, (30), and (31) yield the asserted history-uniform \(o(1)\) expansion without any model-class premise.